% Generated draft: requires mathematical and source review.
\documentclass[11pt]{article}
\usepackage[margin=25mm]{geometry}
\usepackage{fontspec}
\setmainfont{Latin Modern Roman}
\usepackage{amsmath,amssymb,mathtools,unicode-math}
\setmathfont{Latin Modern Math}
\usepackage{enumitem,needspace,xurl,hyperref}
\setlength{\parindent}{0pt}
\setlength{\parskip}{5pt}
\setlength{\emergencystretch}{3em}
\newcommand{\N}{\mathbb{N}}
\newcommand{\Z}{\mathbb{Z}}
\newcommand{\Q}{\mathbb{Q}}
\newcommand{\R}{\mathbb{R}}
\newcommand{\C}{\mathbb{C}}
\newcommand{\F}{\mathbb{F}}
\newcommand{\A}{\mathbb{A}}
\newcommand{\PP}{\mathbb{P}}
\newcommand{\E}{\mathbb{E}}
\newcommand{\eps}{\varepsilon}
\newcommand{\dd}{\,\mathrm{d}}
\begin{document}
\section*{UnsolvedMath ID: 1101221}
\textbf{Source problem code:} catalogue entry
\subsection{Definitions and mathematical statement}
For $n\geq1$, let $F_n$ be the free group of rank $n$ and let
\[
 \operatorname{Out}(F_n)=\operatorname{Aut}(F_n)/\operatorname{Inn}(F_n).
\]
A lattice in a Lie group $L$ is a discrete subgroup $\Gamma$ for which $\Gamma\backslash L$ has finite invariant volume; it is uniform if that quotient is compact. Real rank is the dimension of a maximal real split torus. In the usual finite-centre, no-compact-factor semisimple setting, irreducibility means that the lattice does not decompose across a nontrivial product of the noncompact factors; equivalently its projection to every proper nontrivial collection of simple factors is dense, after the standard finite central quotient.

The source question is: if $\Gamma$ is an irreducible uniform lattice in a higher-rank connected semisimple Lie group $L$, must every homomorphism $\rho:\Gamma\to\operatorname{Out}(F_n)$ have finite image? With higher rank interpreted as $\operatorname{rank}_{\R}L\geq2$, the target quantifiers are
\[
 \forall n\geq1\ \forall\rho\in\operatorname{Hom}(\Gamma,\operatorname{Out}(F_n)),
 \qquad |\rho(\Gamma)|<\infty.
\]
The source does not explicitly state finite-centre or no-compact-factor conventions for $L$; those conventions in the irreducibility explanation are not silently imposed as a correction to the source.

The supplied assessment reports the same conclusion for every irreducible lattice, uniform or nonuniform, in the stated higher-rank setting. It also reports a stronger algebraic criterion: a group $\Gamma$ such that no finite-index subgroup admits a normal subgroup surjecting to $\Z$ has only finite images in $\operatorname{Out}(F_n)$. Explicitly the stated condition is
\[
 \forall\Gamma_0\leq\Gamma\text{ of finite index}\ \forall N\trianglelefteq\Gamma_0,
 \qquad \nexists\text{ surjection }N\longrightarrow\Z.
\]
These are attributed resolution statements; complete theorem hypotheses must be checked against the supplied resolving reference.
\subsection{Short English statement}
Must an irreducible cocompact higher-rank lattice have only finite images in the outer automorphism group of a finite-rank free group? The supplied assessment answers yes, also for noncocompact lattices, citing Bridson--Wade, Corollary B. It attributes the broader normal-subgroup criterion to their Theorem A, and says this also covers hereditarily just-infinite groups that are not virtually cyclic and certain lattices in products of locally compact groups via Bader--Shalom. Those examples and the strengthened scope are reported claims, not blanket assertions about every product lattice. A hereditarily just-infinite group is an infinite group every finite-index subgroup of which has all nontrivial normal subgroups of finite index.
\subsection{Sources}
\begin{itemize}
\item UnsolvedMath permanent ID 1101221, source code catalogue entry, ``Questions in Geometric Group Theory --- Q 12.21.'' Registry curation and metadata: Ulam AI and UnsolvedMath Contributors, CC BY 4.0; updated 2026-08-06.
\item Bestvina, \emph{Questions in Geometric Group Theory} (source list dated 2004), Question 12.21, PDF page 23: \url{https://www.math.utah.edu/~bestvina/eprints/questions-updated.pdf}. Supplied manual transcription accessed 2026-07-29. Original metadata presents an open item without independent current-status verification; status/release and rights review are NEEDS\_REVIEW.
\item Bridson--Wade, \emph{Compositio Mathematica} \textbf{147} (2011), 1573--1580, Theorem A and Corollary B, as cited in the supplied resolving assessment. No article title or locator is supplied.
\item Bader--Shalom, named in the assessment for applications to lattices in products of locally compact groups; no complete citation or exact application hypotheses are provided.
\end{itemize}
\end{document}
